ผู้จัดทำโครงงานขอขอบพระคุณ ผู้ช่วยศาสตราจารย์ ดร.ไอศูรย์ กาญจนสุรัตน์ อาจารย์ที่ปรึกษาโครงงานหลัก ที่ได้กรุณาสละเวลาให้คำแนะนำ คำปรึกษา ถ่ายทอดองค์ความรู้ และชี้แนะแนวทางการดำเนินงานวิจัยอันเป็นประโยชน์อย่างยิ่งตลอดระยะเวลาการดำเนินโครงงาน ตลอดจนช่วยตรวจแก้ไขรายงานโครงงานฉบับนี้ให้มีความสมบูรณ์ถูกต้อง

ขอขอบพระคุณ ผู้ช่วยศาสตราจารย์ นายแพทย์ชัยวัชร์ อภิวาทนสิริ อาจารย์ที่ปรึกษาร่วม สาขาวิชาพยาธิวิทยา คณะแพทยศาสตร์ มหาวิทยาลัยขอนแก่น ที่ได้กรุณาให้คำปรึกษาคำศัพท์ทางการแพทย์เฉพาะทาง และการออกแบบโครงสร้างฟิลด์ข้อมูลตามเกณฑ์มาตรฐานของวิทยาลัยพยาธิแพทย์ (CAP) ซึ่งเป็นพื้นฐานสำคัญที่ทำให้ระบบสามารถตอบโจทย์การใช้งานจริงในบริบททางการแพทย์ได้อย่างมีประสิทธิภาพ

ขอขอบพระคุณคณาจารย์ประจำหลักสูตรปัญญาประดิษฐ์และวิทยาลัยการคอมพิวเตอร์ มหาวิทยาลัยขอนแก่น ทุกท่าน ที่ได้ประสิทธิ์ประสาทวิชาความรู้และทักษะทางวิชาการอันเป็นรากฐานสำคัญในการศึกษาค้นคว้า ตลอดจนขอขอบคุณเพื่อนร่วมชั้นเรียนและผู้มีส่วนเกี่ยวข้องทุกท่านที่ได้ให้ความช่วยเหลือ ร่วมแลกเปลี่ยนความคิดเห็น และเป็นกำลังใจในการดำเนินงานร่วมกันมาโดยตลอด

ท้ายที่สุดนี้ ผู้จัดทำขอกราบขอบพระคุณบิดา มารดา และสมาชิกในครอบครัวทุกท่าน ผู้ซึ่งเป็นแรงผลักดันและคอยให้การสนับสนุนในทุกด้านด้วยความรักและความเข้าใจเสมอมา จนกระทั่งโครงงานฉบับนี้สำเร็จลุล่วงได้ด้วยดี

